\begin{lemma}
\label{lemma-serre-subcategory-is-kernel}
Let $\mathcal{A}$ be an abelian category.
Let $\mathcal{C} \subset \mathcal{A}$ be a Serre subcategory.
There exists an abelian category $\mathcal{A}/\mathcal{C}$
and an exact functor
$$
F : \mathcal{A} \longrightarrow \mathcal{A}/\mathcal{C}
$$
which is essentially surjective and whose kernel is $\mathcal{C}$.
The category $\mathcal{A}/\mathcal{C}$ and the functor $F$ are
characterized by the following universal property: For any exact
functor $G : \mathcal{A} \to \mathcal{B}$ such that
$\mathcal{C} \subset \Ker(G)$ there exists a factorization
$G = H \circ F$ for a unique exact functor
$H : \mathcal{A}/\mathcal{C} \to \mathcal{B}$.
\end{lemma}

\begin{proof}
Consider the set of arrows of $\mathcal{A}$ defined by
the following formula
$$
S = \{f \in \text{Arrows}(\mathcal{A}) \mid
\Ker(f), \Coker(f) \in \Ob(\mathcal{C}) \}.
$$
We claim that $S$ is a multiplicative system. To prove this we have
to check MS1, MS2, MS3, see
Categories, Definition \ref{categories-definition-multiplicative-system}.

\medskip\noindent
It is clear that identities are elements of $S$. Suppose that
$f : A \to B$ and $g : B \to C$ are elements of $S$.
There are exact sequences
$$
\begin{matrix}
0 \to \Ker(f) \to \Ker(gf) \to \Ker(g) \\
\Coker(f) \to \Coker(gf) \to \Coker(g) \to 0
\end{matrix}
$$
Hence it follows that $gf \in S$. This proves MS1. (In fact, a similar
argument will show that $S$ is a saturated multiplicative system, see
Categories, Definition
\ref{categories-definition-saturated-multiplicative-system}.)

\medskip\noindent
Consider a solid diagram
$$
\xymatrix{
A \ar[d]_t \ar[r]_g & B \ar@{..>}[d]^s \\
C \ar@{..>}[r]^f & C \amalg_A B
}
$$
with $t \in S$. Set
$W = C \amalg_A B =  \Coker((t, -g) : A \to C \oplus B)$.
Then $\Ker(t) \to \Ker(s)$ is surjective and
$\Coker(t) \to \Coker(s)$ is an isomorphism. Hence
$s$ is an element of $S$. This proves LMS2 and the proof of RMS2 is dual.

\medskip\noindent
Finally, consider morphisms $f, g : B \to C$ and a morphism $s : A \to B$
in $S$ such that $f \circ s = g \circ s$. This means that
$(f - g) \circ s = 0$. In turn this means that
$I = \Im(f - g) \subset C$ is a quotient of $\Coker(s)$
hence an object of $\mathcal{C}$. Thus $t : C \to C' = C/I$ is an
element of $S$ such that $t \circ (f - g) = 0$, i.e., such that
$t \circ f = t \circ g$. This proves LMS3 and the proof of
RMS3 is dual.

\medskip\noindent
Having proved that $S$ is a multiplicative system we set
$\mathcal{A}/\mathcal{C} = S^{-1}\mathcal{A}$, and we set
$F$ equal to the localization functor $Q$. By
Lemma \ref{lemma-localization-abelian}
the category $\mathcal{A}/\mathcal{C}$ is abelian and $F$ is exact.
If $X$ is in the kernel of $F = Q$, then by
Lemma \ref{lemma-kernel-localization}
we see that $0 : X \to Z$ is an element of $S$ and hence
$X$ is an object of $\mathcal{C}$, i.e., the kernel of
$F$ is $\mathcal{C}$.
Finally, if $G$ is as in the statement of the lemma, then $G$ turns
every element of $S$ into an isomorphism. Hence we obtain the
functor $H : \mathcal{A}/\mathcal{C} \to \mathcal{B}$ from
the universal property of localization, see
Categories, Lemma \ref{categories-lemma-properties-left-localization}.
We still have to show the functor $H$ is exact.
To do this it suffices to show that $H$ commutes
with taking kernels and cokernels, see Lemma \ref{lemma-exact-functor}.
Let $A \to B$ be a morphism in $\mathcal{A}/\mathcal{C}$.
We may represent $A \to B$ as $fs^{-1}$ where $s : A' \to A$
is in $S$ and $f : A' \to B$ an arbitrary morphism of $\mathcal{A}$.
Since $F = Q$ maps $s$ to an isomorphism in the quotient category
$\mathcal{A}/\mathcal{C}$, it suffices to show that $H$ commutes with taking
kernels and cokernels of morphisms $f : A \to B$ of $\mathcal{A}$.
But here we have $H(f) = G(f)$ and the result follows
from the fact that $G$ is exact.
\end{proof}
